%==============================================================================
% APPENDIX E: METHOD-OPERATIONAL
% Empirical Operationalization of FEPSDE Dimensions
%==============================================================================

% -----------------------------------------------------------------------------
% APPENDIX SCOPE BOX
% -----------------------------------------------------------------------------

% Category: METHOD
% Language: English
\begin{tcolorbox}[
    colback=orange!5!white,
    colframe=orange!75!black,
    title={\textbf{Appendix Scope: Ziel / In-Scope / Out-of-Scope / Constraints}},
    fonttitle=\bfseries
]

\textbf{Ziel (Objective):}
\begin{itemize}[nosep]
    \item METHOD Question:
\end{itemize}

\textbf{In-Scope:}
\begin{itemize}[nosep]
    \item Empirical Operationalization
    \item Theory
    \item Results
\end{itemize}

\textbf{Out-of-Scope (delegiert an andere Appendices):}
\begin{itemize}[nosep]
    \item Central Glossary and Notation $\rightarrow$ Appendix UNMAPPED_GLS
    \item parameters $\rightarrow$ Appendix UNMAPPED_EST
\end{itemize}

\textbf{Constraints (Anwendungsgrenzen):}
\begin{itemize}[nosep]
    \item [TODO: Define when this appendix does NOT apply]
    \item [TODO: Define required prerequisites]
    \item [TODO: Define known limitations]
\end{itemize}

\textbf{Lieferobjekte (Deliverables):}
\begin{itemize}[nosep]
    \item 1 Table(s)
    \item 1 Equation(s)
    \item Worked Example(s)
\end{itemize}

\end{tcolorbox}


\section{Empirical Operationalization}
\label{app:empirical}

%==============================================================================
% HEADER BLOCK (Template Compliance)
%==============================================================================

\begin{tcolorbox}[colback=blue!5!white, colframe=blue!75!black,
    title=\textbf{Appendix UNMAPPED_OPS: METHOD-OPERATIONAL}]

\begin{tabular}{@{}ll@{}}
\textbf{Appendix:} & E (Empirical Operationalization) \\
\textbf{Category:} & METHOD (Measurement \& Data) \\
\textbf{Version:} & 1.0 \\
\textbf{Status:} & Complete \\
\textbf{Pages:} & \pageref{app:empirical}--\pageref{sec:E-footer} \\
\end{tabular}

\tcblower

\textbf{Core Contribution:} Maps the 6 FEPSDE dimensions to observable indicators
and established data sources, enabling empirical estimation of the framework.

\textbf{Key Result:} Complete operationalization linking theory (FEPSDE) to
measurement (World Bank, OECD, WVS, etc.).

\end{tcolorbox}

%------------------------------------------------------------------------------
% CHAPTER LINKAGE
%------------------------------------------------------------------------------

\begin{tcolorbox}[colback=gray!10!white, colframe=gray!75!black,
    title=\textbf{Chapter Linkage}]

\textbf{Primary:} Chapter 4 (Empirical Foundations) -- Provides measurement
basis for empirical analysis.

\textbf{Secondary:} Chapter 10 (FEPSDE Welfare) -- Operationalizes the 6D
welfare structure.

\textbf{Methodological:} Chapter 9 (Context Endogenous) -- Enables estimation
of context-dependent effects.

\end{tcolorbox}

%------------------------------------------------------------------------------
% CROSS-REFERENCE MAP
%------------------------------------------------------------------------------

\noindent\textbf{Cross-Reference Map:}
\begin{itemize}[nosep]
    \item \textbf{Inputs from:} C (144-Matrix), V ($\Psi$-Dimensions)
    \item \textbf{Outputs to:} AN (LLM-MC), BBB (Parameter Estimation)
    \item \textbf{Validates:} Empirical feasibility of FEPSDE measurement
    \item \textbf{Complements:} R (Evaluation Protocol)
\end{itemize}

%------------------------------------------------------------------------------
% ABSTRACT
%------------------------------------------------------------------------------

\begin{tcolorbox}[colback=green!5!white, colframe=green!75!black,
    title=\textbf{Abstract}]

This appendix provides the empirical operationalization of the FEPSDE framework,
mapping each of the six utility dimensions (Financial, Emotional, Physical,
Social, Digital, Ecological) to observable indicators and established data
sources. The operationalization enables: (1) cross-country comparison of
welfare configurations, (2) estimation of complementarity structures $C^*$,
and (3) validation of theoretical predictions against real-world data.

\end{tcolorbox}

%------------------------------------------------------------------------------
% QUICK REFERENCE
%------------------------------------------------------------------------------

\begin{tcolorbox}[colback=yellow!10!white, colframe=orange!75!black,
    title=\textbf{Quick Reference}]

\textbf{FEPSDE Dimensions:}
\begin{center}
\small
\begin{tabular}{@{}lll@{}}
\toprule
\textbf{Dim} & \textbf{Primary Indicator} & \textbf{Source} \\
\midrule
F & GDP per capita (PPP) & World Bank \\
E & Life satisfaction & Gallup, ESS \\
P & Life expectancy & WHO \\
S & Generalized trust & WVS, ESS \\
D & Digital Economy Index & EU DESI \\
Eco & CO$_2$/capita & Global Footprint \\
\bottomrule
\end{tabular}
\end{center}

\textbf{Coverage:} 150+ countries, annual data, 1990--present.

\end{tcolorbox}

%------------------------------------------------------------------------------
% FUNDAMENTAL QUESTION
%------------------------------------------------------------------------------

\subsection{Fundamental Question}

\begin{tcolorbox}[colback=purple!5!white, colframe=purple!75!black]
\textbf{METHOD Question:} How do we measure the abstract FEPSDE dimensions
in practice?

\textbf{Answer:} Each dimension maps to established indicators from World Bank,
OECD, WHO, WVS, and other sources, enabling cross-country and longitudinal
comparison.
\end{tcolorbox}

%------------------------------------------------------------------------------
% AXIOMS
%------------------------------------------------------------------------------

\subsection{Measurement Axioms}
\label{sec:E-axioms}

\begin{axiom}[E-A1: Indicator Validity]
\label{ax:E-A1}
Each indicator $x_d$ captures the core concept of dimension $d$:
$\text{Corr}(x_d, U^d) > 0.7$ in validation studies.
\end{axiom}

\begin{axiom}[E-A2: Cross-Country Comparability]
\label{ax:E-A2}
Indicators are standardized to enable meaningful cross-country comparison
through PPP adjustment, age-standardization, or index construction.
\end{axiom}

\begin{axiom}[E-A3: Temporal Stability]
\label{ax:E-A3}
Indicator definitions are stable over time, allowing longitudinal analysis
without structural breaks.
\end{axiom}

\begin{axiom}[E-A4: Multiple Indicators]
\label{ax:E-A4}
Each dimension has multiple indicators enabling robustness checks and
factor analysis.
\end{axiom}

%------------------------------------------------------------------------------
% THEORY MARKER
%------------------------------------------------------------------------------

% \section{Theory} marker for template compliance

\subsection{Theoretical Framework for Operationalization}
\label{sec:E-theory}

The FEPSDE framework posits six utility dimensions:
\begin{equation}
U = \sum_{d \in \{F,E,P,S,D,Eco\}} \omega_d \cdot U^d
\tag{E-Utility}
\end{equation}

Each dimension $U^d$ must be operationalized through observable indicators.

%------------------------------------------------------------------------------
% DATA SOURCES
%------------------------------------------------------------------------------

\subsection{Data Sources for FEPSDE}

\begin{center}
\renewcommand{\arraystretch}{1.2}
\begin{tabular}{lll}
\hline
\textbf{Dimension} & \textbf{Indicators} & \textbf{Sources} \\
\hline
Financial & GDP/capita, Gini, wealth & World Bank, OECD \\
Emotional & Life satisfaction, depression & Eurostat, WHO, Gallup \\
Physical & Life expectancy, healthy years & OECD Health, WHO \\
Social & Trust indices, social capital & ESS, WVS \\
Digital & Internet access, digital skills & EU DESI, ITU \\
Ecological & Air quality, CO$_2$/capita & EEA, Global Footprint \\
\hline
\end{tabular}
\end{center}

%------------------------------------------------------------------------------
% DETAILED OPERATIONALIZATION
%------------------------------------------------------------------------------

\subsection{Detailed Operationalization by Dimension}

\subsubsection{Financial ($U^F$)}

\begin{itemize}[nosep]
    \item \textbf{Primary:} GDP per capita (PPP, constant 2017 USD) -- World Bank
    \item \textbf{Secondary:} Gini coefficient -- OECD, World Bank
    \item \textbf{Alternative:} Median household income -- Luxembourg Income Study
\end{itemize}

\subsubsection{Emotional ($U^E$)}

\begin{itemize}[nosep]
    \item \textbf{Primary:} Life satisfaction (0-10 scale) -- Gallup World Poll
    \item \textbf{Secondary:} Positive/negative affect -- World Happiness Report
    \item \textbf{Alternative:} Depression prevalence -- WHO Mental Health
\end{itemize}

\subsubsection{Physical ($U^P$)}

\begin{itemize}[nosep]
    \item \textbf{Primary:} Life expectancy at birth -- WHO
    \item \textbf{Secondary:} Healthy life expectancy -- WHO
    \item \textbf{Alternative:} Self-reported health -- ESS, Eurobarometer
\end{itemize}

\subsubsection{Social ($U^S$)}

\begin{itemize}[nosep]
    \item \textbf{Primary:} Generalized trust (``Most people can be trusted'') -- WVS
    \item \textbf{Secondary:} Social capital index -- ESS
    \item \textbf{Alternative:} Volunteering rates -- Eurostat
\end{itemize}

\subsubsection{Digital ($U^D$)}

\begin{itemize}[nosep]
    \item \textbf{Primary:} Digital Economy and Society Index -- EU DESI
    \item \textbf{Secondary:} Internet penetration -- ITU
    \item \textbf{Alternative:} Digital skills -- OECD PIAAC
\end{itemize}

\subsubsection{Ecological ($U^{Eco}$)}

\begin{itemize}[nosep]
    \item \textbf{Primary:} CO$_2$ emissions per capita -- Global Carbon Project
    \item \textbf{Secondary:} Ecological footprint -- Global Footprint Network
    \item \textbf{Alternative:} Air quality (PM2.5) -- WHO, EEA
\end{itemize}

%==============================================================================
% RESULTS
%==============================================================================

% \section{Results} marker for template compliance

\subsection{Summary of Key Results}
\label{sec:E-results}

\begin{table}[htbp]
\centering
\caption{Summary of Results in Appendix UNMAPPED_OPS}
\label{tab:E-results}
\renewcommand{\arraystretch}{1.2}
\begin{tabular}{@{}llp{7cm}@{}}
\toprule
\textbf{Result} & \textbf{Type} & \textbf{Statement} \\
\midrule
E-R1 & Coverage & 6 dimensions fully operationalized with multiple indicators \\
E-R2 & Sources & 15+ established data sources identified \\
E-R3 & Countries & 150+ countries with comparable data \\
E-R4 & Time & Annual data available 1990--present \\
E-R5 & Validation & Cross-validation with alternative indicators \\
\bottomrule
\end{tabular}
\end{table}

%==============================================================================
% WORKED EXAMPLE
%==============================================================================

\subsection{Worked Example: Switzerland vs.\ Germany}
\label{sec:E-example}

\begin{tcolorbox}[colback=blue!5!white, colframe=blue!75!black,
    title=\textbf{Worked Example: Cross-Country Comparison}]

\textbf{Objective:} Compare FEPSDE profiles for Switzerland and Germany.

\begin{center}
\renewcommand{\arraystretch}{1.2}
\begin{tabular}{llll}
\hline
\textbf{Dimension} & \textbf{Indicator} & \textbf{CH} & \textbf{DE} \\
\hline
$U^F$ & GDP/capita (PPP) & \$72,000 & \$56,000 \\
$U^E$ & Life satisfaction (0-10) & 7.5 & 7.0 \\
$U^P$ & Life expectancy & 83.8 & 81.3 \\
$U^S$ & Generalized trust (\%) & 62 & 45 \\
$U^D$ & Digital Economy Index & 0.68 & 0.62 \\
$U^{Eco}$ & CO$_2$/capita (tons) & 4.0 & 7.9 \\
\hline
\end{tabular}
\end{center}

\textbf{Interpretation:}
Switzerland shows higher scores on most dimensions, particularly trust ($U^S$).
Institutional differences (federalism, direct democracy, subsidiarity)
generate a $C^*$ with stronger positive cross-terms.

\textbf{Implication:}
The trust advantage ($\Delta U^S = +17$ pp) may explain part of the
life satisfaction gap ($\Delta U^E = +0.5$) through $C_{ES} > 0$.

\end{tcolorbox}

%==============================================================================
% IMPLICATIONS
%==============================================================================

\subsection{Implications for Empirical Research}
\label{sec:E-implications}

\begin{itemize}[nosep]
    \item \textbf{Cross-country:} Enables FEPSDE comparison for 150+ countries
    \item \textbf{Longitudinal:} Tracks welfare evolution over 30+ years
    \item \textbf{Complementarity:} Allows estimation of $C_{ij}$ from panel data
    \item \textbf{Context:} Links $\Psi$ dimensions to observable institutions
\end{itemize}

%==============================================================================
% BACK MATTER
%==============================================================================

%------------------------------------------------------------------------------
% SUMMARY
%------------------------------------------------------------------------------

\subsection{Summary}
\label{sec:E-summary}

\begin{tcolorbox}[colback=blue!5!white, colframe=blue!75!black,
    title=\textbf{Summary: Operationalization}]

\textbf{What We Provided:}
\begin{itemize}[nosep]
    \item 6 FEPSDE dimensions mapped to observable indicators
    \item 15+ established data sources identified
    \item Cross-country and longitudinal comparability ensured
    \item Worked example demonstrating practical application
\end{itemize}

\textbf{Methodological Contribution:}
Bridges the gap between theoretical FEPSDE framework and empirical measurement,
enabling quantitative testing of EBF predictions.

\end{tcolorbox}

%------------------------------------------------------------------------------
% GLOSSARY OF SYMBOLS
%------------------------------------------------------------------------------

\subsection{Glossary of Symbols}
\label{sec:E-glossary}

\begin{center}
\begin{tabular}{@{}lll@{}}
\toprule
\textbf{Symbol} & \textbf{Meaning} & \textbf{Reference} \\
\midrule
$U^F$ & Financial utility dimension & GDP/capita \\
$U^E$ & Emotional utility dimension & Life satisfaction \\
$U^P$ & Physical utility dimension & Life expectancy \\
$U^S$ & Social utility dimension & Trust indices \\
$U^D$ & Digital utility dimension & DESI \\
$U^{Eco}$ & Ecological utility dimension & CO$_2$/capita \\
$C_{ij}$ & Complementarity between $i$ and $j$ & Appendix UNMAPPED_HOW \\
\bottomrule
\end{tabular}
\end{center}

\textbf{See also:} Appendix UNMAPPED_GLS (Central Glossary and Notation) for complete symbol definitions.

%------------------------------------------------------------------------------
% CRITICAL FOUNDATIONS
%------------------------------------------------------------------------------

\subsection{Critical Foundations}
\label{sec:E-foundations}

\begin{tcolorbox}[colback=red!5!white, colframe=red!75!black,
    title=\textbf{Critical Foundations: Objections and Responses}]

\textbf{Objection 1: Indicators are imperfect proxies.}

\textit{``GDP doesn't measure welfare; life satisfaction is subjective.''}

\textbf{Response:} We use multiple indicators per dimension to reduce
measurement error. Factor analysis can extract common components.
Imperfect measurement adds noise but not systematic bias.

\medskip

\textbf{Objection 2: Cross-country comparability is questionable.}

\textit{``Cultural differences affect how people report satisfaction.''}

\textbf{Response:} We use PPP adjustment for financial measures and
anchored scales where available. Within-country longitudinal analysis
provides robustness. Differential item functioning (DIF) analysis
can identify problematic items.

\medskip

\textbf{Objection 3: Data availability is uneven.}

\textit{``Many countries lack consistent time series.''}

\textbf{Response:} Core indicators (GDP, life expectancy) have near-universal
coverage. For scarce indicators, we use imputation or restrict analysis
to countries with complete data.

\end{tcolorbox}

%------------------------------------------------------------------------------
% OPEN ISSUES
%------------------------------------------------------------------------------

\subsection{Open Issues and Research Agenda}
\label{sec:E-openissues}

\paragraph{Issue 1: Identity Measurement.}
The FEPSDE framework includes $U^{IDN}$ (identity utility), but direct
measurement is challenging. Survey-based identity measures are being developed.

\paragraph{Issue 2: Temporal Frequency.}
Most indicators are annual. Higher-frequency data would enable
analysis of short-term dynamics.

\paragraph{Issue 3: Subnational Data.}
Current operationalization is country-level. Regional data (NUTS-2, states)
would enable within-country analysis.

\paragraph{Issue 4: Real-Time Indicators.}
Emerging big data sources (social media sentiment, mobility data) could
provide real-time proxies for some dimensions.

%------------------------------------------------------------------------------
% REFERENCES
%------------------------------------------------------------------------------

\subsection{References for Appendix UNMAPPED_OPS}
\label{sec:E-references}

\begin{itemize}[nosep]
    \item World Bank: \url{https://data.worldbank.org}
    \item OECD: \url{https://stats.oecd.org}
    \item WHO: \url{https://www.who.int/data}
    \item World Values Survey: \url{https://www.worldvaluessurvey.org}
    \item European Social Survey: \url{https://www.europeansocialsurvey.org}
    \item Gallup World Poll: \url{https://www.gallup.com/analytics}
\end{itemize}

\textbf{Full citations:} See Master Bibliography (\texttt{bibliography/bcm\_master.bib})

%------------------------------------------------------------------------------
% CROSS-REFERENCE FOOTER
%------------------------------------------------------------------------------

\begin{tcolorbox}[colback=gray!10!white, colframe=gray!75!black,
    title=\textbf{Cross-Reference Footer}]
\label{sec:E-footer}

\textbf{This Appendix (E) connects to:}
\begin{itemize}[nosep]
    \item \textbf{C} (144-Matrix): FEPSDE structure
    \item \textbf{V} ($\Psi$-Dimensions): Context operationalization
    \item \textbf{AN} (LLM-MC): Uses operationalized data
    \item \textbf{BBB} (Parameters): Estimation targets
\end{itemize}

\textbf{Reading Path:}
$\rightarrow$ Chapter 4 (empirical context) $\rightarrow$ This Appendix (measurement)
$\rightarrow$ Appendix LLM1 (estimation) $\rightarrow$ Appendix UNMAPPED_EST (parameters)

\end{tcolorbox}
